\chapter{The Topology of Gold's Bands}
\label{ch:topology}

Everything computed so far --- energies, bands, densities of states --- is a
\emph{number that varies continuously} when you change the lattice constant
or the functional a little. This chapter extracts a quantity of a different
kind: an \emph{integer}, $0$ or $1$, that is a property of the
wavefunctions rather than of the energies and that cannot change under
small perturbations. It is the $\mathbb{Z}_2$ topological invariant.
Gold's value is $1$ (``nontrivial''), and the well-known surface state of
the Au(111) surface is its experimental fingerprint.

Nothing new has to be built. The three ingredients are already in your
hands: the even/odd character of zone-boundary states
(Section~\ref{sec:nfe}), parity and time reversal as permutations of the
coefficient array (Section~\ref{sec:symmetry}), and the fixed-potential
band solver \textsc{SolveBands} (Section~\ref{sec:nscf}). The chapter
has two halves. Sections~\ref{sec:berry}--\ref{sec:z2def} explain what the
invariant \emph{is}, in terms of where the electronic charge sits in the
unit cell and how it moves; this leads to a general algorithm, the Wilson
loop. Sections~\ref{sec:fukane}--\ref{sec:goldparity} then show that
inversion symmetry collapses the whole computation to multiplying eight
signs. You will do both, and they must agree.

\section{What ``topological'' means}

A doughnut and a coffee mug both have exactly one hole. You can deform one
into the other smoothly; you cannot get rid of the hole without tearing.
The number of holes is a \emph{topological invariant}: an integer that
survives every smooth deformation.

Band structures have invariants of the same kind. Take a group of bands
separated from all higher bands by an energy gap. Deform the Hamiltonian
smoothly (strain the crystal, change the potential) without ever closing
that gap and without breaking time-reversal symmetry. Some properties of
the group of wavefunctions cannot change under such a deformation, and for
systems with spin--orbit coupling one of them is a single bit
$\nu_0 \in \{0, 1\}$.
\begin{itemize}[leftmargin=2em]
\item $\nu_0 = 0$: the bands can be deformed into those of isolated atoms
pulled far apart. Nothing special.
\item $\nu_0 = 1$: they cannot. Somewhere in the Brillouin zone the bands
are ordered ``the wrong way round'' compared with separated atoms --- a
\textbf{band inversion} --- and no smooth change can undo it.
\end{itemize}
Why care about a bit? Because of what happens at a surface. Outside the
crystal is vacuum, which is trivially $\nu_0 = 0$. If the crystal has
$\nu_0 = 1$, the bit has to change across the surface, and it can only
change where the gap closes. So the surface is \emph{forced} to host
electronic states inside the bulk gap, and SOC makes them spin-polarized
with the spin locked to the direction of motion. Those surface states are
what experiments see.

To say precisely what the bit measures we need one concept that has not
appeared yet.

\section{Where is the charge? Berry phases and Wannier centers}
\label{sec:berry}

\paragraph{The question.} A filled band contributes a lattice-periodic
charge density. Ask a seemingly innocent question: \emph{where, within the
unit cell, is the center of that band's charge?} The density alone cannot
answer it --- any periodic density can be assigned to cells in many ways.
The answer is hidden in how the cell-periodic functions $u_{n\kk}$
\emph{change with} $\kk$.

\paragraph{One band in one dimension.} Take a single isolated band on a 1D
lattice of period $a$, with $u_k(x)$ on a mesh $k_j = 2\pi j/(N a)$,
$j = 0,\dots,N{-}1$. Each $u_{k_j}$ comes out of the eigensolver with an
arbitrary phase, so an individual overlap $\braket{u_{k_j}|u_{k_{j+1}}}$
means nothing by itself. But multiply the overlaps around the closed loop
of the Brillouin zone,
\begin{equation}
\mathcal W = \braket{u_{k_0}|u_{k_1}}\braket{u_{k_1}|u_{k_2}}\cdots
\braket{u_{k_{N-1}}|u_{k_N}},
\qquad u_{k_N}(x) \equiv \ee^{-\ii 2\pi x/a}\,u_{k_0}(x),
\label{eq:berry1d}
\end{equation}
and every state appears once as a bra and once as a ket: all the arbitrary
phases cancel. (The last factor closes the loop: $k_N = k_0 + 2\pi/a$ is
the same Bloch state as $k_0$, and $\psi = \ee^{\ii kx}u_k$ then fixes
$u_{k_N}$ as stated.) The phase of this product,
\begin{equation}
\phi = -\,\mathrm{Im}\ln \mathcal W \pmod{2\pi},
\end{equation}
is the \textbf{Berry phase} of the band (in one dimension also called the
Zak phase). It is a genuine, computable property of the band, and its
meaning is the answer to our question:
\begin{equation}
\boxed{\;\bar x = \frac{\phi}{2\pi}\,a \pmod a\;}
\label{eq:wcc}
\end{equation}
is the center of charge of the band, measured from the origin of the cell
--- the \textbf{Wannier charge center}. (``Wannier'' because it is the
mean position of the localized Wannier function built from the band;
``mod $a$'' because the charge center is only defined up to a lattice
translation.)

\emph{Why this is plausible.} In $\kk$-space the position operator acts as
$\ii\,\partial/\partial k$ on $u_k$ (just as momentum acts as
$-\ii\,\partial/\partial x$ in real space). For neighbouring mesh points,
$\braket{u_k|u_{k+\delta}} \approx 1 + \delta\braket{u_k|\partial_k u_k}
\approx \exp\!\big(-\ii\delta\, \braket{u_k|\ii\partial_k|u_k}\big)$, so
$-\mathrm{Im}\ln$ of the product adds up
$\delta \times(\text{local expectation of position})$ around the zone;
dividing by the length $2\pi/a$ of the zone gives the average
position. A check you can do by hand: for a band made of one atomic
orbital sitting at position $x_0$ in each cell (tight-binding limit),
$u_k(x) \propto \sum_R \ee^{-\ii k(x - x_0 - R)}\varphi(x-x_0-R)$, and
Eq.~\eqref{eq:berry1d} gives $\phi = 2\pi x_0/a$ (all interior overlaps
are real and positive on a fine mesh; the whole phase comes from the
closing factor): $\bar x = x_0$. The charge is where the orbital is.

\paragraph{Several bands: the Wilson loop.} For a group of $N_{\mathrm{occ}}$
bands (which may cross or be degenerate among themselves, as long as the
group is separated from the bands above), individual bands have no
identity, so the overlaps become matrices
\begin{equation}
M^{(\kk,\kk')}_{mn} = \braket{u_{m\kk}|u_{n\kk'}}
= \sum_{\GG,s} c^{s\,*}_{m\kk}(\GG)\, c^{s}_{n\kk'}(\GG),
\qquad m,n = 1,\dots,N_{\mathrm{occ}},
\label{eq:overlapM}
\end{equation}
(sum over the plane waves $\GG$ common to both spheres), and the loop
product becomes a matrix, the \textbf{Wilson loop}
\begin{equation}
\mathcal W = M^{(\kk_0,\kk_1)}\, M^{(\kk_1,\kk_2)} \cdots
M^{(\kk_{N-1},\,\kk_N)} .
\label{eq:wilson}
\end{equation}
Under a change of basis within the group at any intermediate $\kk_j$ the
product is unchanged; a change at $\kk_0$ conjugates $\mathcal W$ by a
unitary matrix. Its \emph{eigenvalues} are therefore well defined. On a
fine mesh they have modulus $\approx 1$, and their phases
$\theta_j \in (-\pi,\pi]$, $j = 1,\dots,N_{\mathrm{occ}}$, give
$N_{\mathrm{occ}}$ charge centers $\bar x_j = a\,\theta_j/2\pi$: the
positions of the $N_{\mathrm{occ}}$ electrons per cell that the group of
bands holds. (Sign conventions for $\theta$ differ between texts; none of
what follows depends on the choice.)

\paragraph{In three dimensions: centers that depend on $\kk_\perp$.} In
the crystal, run the loop along one reciprocal direction only: fix the
fractional coordinates $k_2, k_3$ and let $k_1$ go once around,
$\kk_j = (j/N_1,\,k_2,\,k_3)$. The last overlap again closes the loop with
the \emph{same} states as at $\kk_0$, relabeled: since
$u_{\kk+\mathbf b_1}(\rr) = \ee^{-\ii\mathbf b_1\cdot\rr}u_{\kk}(\rr)$,
\begin{equation}
c_{\kk_0+\mathbf b_1}(\GG) = c_{\kk_0}(\GG + \mathbf b_1)
\label{eq:closure}
\end{equation}
--- a shift of the Miller index $m_1$ by one; do \emph{not} solve again at
$\kk_0 + \mathbf b_1$, which would return the same states with new random
phases and destroy the cancellation. The result is a set of centers
$\bar x_j(k_2,k_3)$ along $\mathbf a_1$ that depend on the transverse
momentum: for each $(k_2,k_3)$ the crystal behaves like a 1D chain along
$\mathbf a_1$, and $\theta_j(k_2,k_3)$ tells you where its electrons sit.

\section{The \texorpdfstring{$\mathbb{Z}_2$}{Z2} invariant: how the centers flow}
\label{sec:z2def}

Fix $k_3$ and watch the $N_{\mathrm{occ}}$ centers $\theta_j(k_2)$ as $k_2$
goes across the zone. They trace continuous curves on a cylinder ($\theta$
is an angle). Time-reversal symmetry constrains these curves in two ways.
\begin{itemize}[leftmargin=2em]
\item $\theta_j(-k_2) = \theta_j(k_2)$ as a set: the picture for
$k_2\in[-\tfrac12,0]$ is the mirror image of that for $[0,\tfrac12]$, so
half the zone suffices.
\item At $k_2 = 0$ and $k_2 = \tfrac12$ the 1D chain is itself
time-reversal invariant, and by Kramers' theorem its centers are
\emph{doubly degenerate}: the curves meet in pairs at both ends of the
half zone.
\end{itemize}
Between the two ends the pairs split. The only question topology asks is
how they recombine (compare the two panels of Fig.~\ref{fig:wcc}): each curve can return to the
partner it started with, or the curves can \emph{switch partners}, each
ending up paired with a curve that started in a different pair. Switching
partners cannot be undone by any smooth deformation --- to do so two curves
would have to pass through each other, and Kramers degeneracy protects the
crossings at the ends. The parity of the number of partner switches is the
invariant of the plane:
\begin{equation}
z(k_3) = \Big(\text{number of crossings of the line }
\theta = \theta_{\mathrm{ref}} \text{ by the curves } \theta_j(k_2),\
0\le k_2\le\tfrac12\Big) \bmod 2 .
\label{eq:zcount}
\end{equation}
If every curve returns to its partner, any horizontal line is crossed an
even number of times; if partners are switched, some curve must travel
from one pair to the next and crosses the line an odd number of times.
The result does not depend on $\theta_{\mathrm{ref}}$. Physically: as $k_2$
advances through half the zone, a trivial band group shuffles its charge
centers and puts them back; a nontrivial one hands a charge center over
from one Kramers pair to the next, a net flow that cannot be stopped
without closing the gap.

A three-dimensional crystal has two independent time-reversal-invariant
planes of this kind for each direction, $k_3 = 0$ and $k_3 = \tfrac12$.
The \textbf{strong index} and one of the three \textbf{weak indices} are
\begin{equation}
\nu_0 = \big[z(k_3{=}0) + z(k_3{=}\tfrac12)\big] \bmod 2,
\qquad
\nu_3 = z(k_3{=}\tfrac12),
\label{eq:nu0fromz}
\end{equation}
and $\nu_1$, $\nu_2$ follow from the planes $k_1 = \tfrac12$,
$k_2 = \tfrac12$. The strong index is the robust one: $\nu_0 = 1$ means
the two parallel planes disagree, which is what guarantees a protected
state on \emph{every} surface (in a metal such as gold: inside the
projected direct gap, not in an absolute gap;
Section~\ref{sec:directgap}). The full label is written
$(\nu_0;\nu_1\nu_2\nu_3)$.

\section{But gold is a metal}
\label{sec:directgap}

The construction needs a group of bands separated from the rest by a gap,
and a metal has no gap at the Fermi level. Look closer, though: the
definition never mentions the Fermi level or the occupations. It needs
only that the highest band of the group and the next band do not touch
\emph{at any single $\kk$} --- a ``direct gap'' everywhere, even if the
top of the group at one $\kk$ lies above the bottom of the next band at
another. Gold has exactly this between Kramers pair 6 and pair 7 (spinor
bands 12 and 13). Pairs 1--6 are the five $d$ bands plus the $sp$ band
that crosses $E_F$; pair 7 lies above them at every $\kk$:
by $5.1$\,eV at $X$, $4.3$\,eV at $L$ and $19$\,eV at $\Gamma$ in the
reference calculation of Chapter~\ref{ch:verify}. Away from the TRIM the
gap is much smaller: it falls to $0.07$\,eV on the line $W$--$L$, and a
search over the zone found a minimum of about $0.02$\,eV at a point on a
mirror plane, $\kk \approx \tfrac{2\pi}{a}(0.68, 0.68, 0.18)$ (equivalent,
by a reciprocal lattice vector and cubic symmetry, to
$\tfrac{2\pi}{a}(0.82, 0.32, 0.32)$ inside the first zone). Small, but
not zero.
So the group is $N_{\mathrm{occ}} = 12$ spinor bands, six Kramers pairs,
regardless of the fact that part of pair 6 is empty. (That gold's
invariant defined this way is nontrivial was pointed out by Yan and
co-workers in 2015.) The first step of any topology calculation on a metal
is therefore to \emph{verify the direct gap} on a fine mesh.

Two practical warnings follow from these numbers. First, a coarse mesh
does not see the small gaps: the $8^3$ SCF mesh reports a minimum of
$0.54$\,eV and a $20^3$ mesh $0.27$\,eV, so finish the check with a local
minimization of $\varepsilon_{13}-\varepsilon_{12}$ started from the
smallest values found. Second, it is spin--orbit coupling that keeps the
two groups apart at such points: with $k^\ell = 0$ the same search finds
the sixth and seventh bands \emph{touching} (gap below $10^{-4}$\,eV) at
several points of the zone --- without SOC, bands of different mirror
symmetry are free to cross. The invariant of gold is a spin--orbit effect not because
SOC changes any parity, but because SOC makes the group of bands well
defined.

\section{The shortcut for inversion-symmetric crystals: parities}
\label{sec:fukane}

Following charge centers across the zone works for any crystal. For a
crystal with inversion symmetry Fu and Kane (2007) proved that the answer
is already fixed by the parities $\xi_m(\Gamma_i) = \pm1$ of the Kramers
pairs at the eight TRIM $\Gamma_i$ (Section~\ref{sec:symmetry}). Multiply
the parities of the $N$ pairs of the group at each TRIM, then multiply the
eight results:
\begin{equation}
\delta_i = \prod_{m=1}^{N} \xi_m(\Gamma_i) ,
\qquad\qquad
\boxed{\;(-1)^{\nu_0} = \prod_{i=1}^{8} \delta_i\;}
\label{eq:fukane}
\end{equation}
The weak index $\nu_j$ is the same product restricted to the four TRIM
with $n_j = 1$.

\paragraph{Why parities know about charge centers.} The connection to the
previous section is short enough to sketch. Consider a Wilson loop along
$k_1$ at one of the special values of $(k_2,k_3)$ where the 1D chain is
time-reversal invariant; the loop passes through exactly two TRIM,
$\Gamma_a$ (at $k_1 = 0$) and $\Gamma_b$ (at $k_1=\tfrac12$). Inversion
maps a charge center $\bar x$ to $-\bar x$, so a center that is alone must
sit at a point that is its own mirror image: $\bar x = 0$ (on the atom,
$\theta = 0$) or $\bar x = a/2$ (between atoms, $\theta = \pi$). Which one?
Exactly as in the two-plane-wave model of Section~\ref{sec:nfe}: a band
whose states at $\Gamma_a$ and $\Gamma_b$ have the \emph{same} parity is
centered on the atom; if the parities are \emph{opposite}, on the midpoint.
For the whole group (centers may also come in mirror-image sets
$\pm\theta$, which cancel),
\begin{equation}
\prod_{\text{pairs } m} \ee^{\ii\theta_m}
= \prod_{m=1}^{N} \xi_m(\Gamma_a)\,\xi_m(\Gamma_b) = \delta_a\delta_b ,
\label{eq:zakparity}
\end{equation}
one factor per Kramers pair: the number of pairs centered at
$\theta = \pi$ is odd precisely when $\delta_a\delta_b = -1$. Now compare
the two ends of the half zone in Eq.~\eqref{eq:zcount}. If the number of
pairs sitting at $\theta=\pi$ changes from even at $k_2 = 0$ to odd at
$k_2=\tfrac12$, one center must have travelled from $0$ to $\pi$ and
thereby crossed a reference line once: $z = 1$. Hence
$(-1)^{z} = \delta_a\delta_b\delta_c\delta_d$ over the four TRIM of the
plane, and Eq.~\eqref{eq:nu0fromz} multiplies the two planes together into
Eq.~\eqref{eq:fukane}. You will verify Eq.~\eqref{eq:zakparity}
numerically below.

\paragraph{The intuition.} For atoms pulled far apart, every band is built
from one kind of atomic orbital and has the same parity $(-1)^\ell$ at
every TRIM, so each band contributes $(\pm1)^8 = +1$ and $\nu_0 = 0$. To
get $\nu_0 = 1$, an even state and an odd state must have swapped places,
relative to the atomic ordering, at an odd number of TRIM. The formula
counts band inversions.

\section{Algorithm 1: the invariant from parities}
\label{sec:goldparity}

Your eigensolver returns \emph{some} orthonormal basis of each degenerate
level, not parity eigenstates. That does not matter. With the permutation
of Eq.~\eqref{eq:parityperm}, form the matrix of $\hat P$ among the lowest
$2N$ spinor bands at a TRIM,
\begin{equation}
\Pi_{mn} = \braket{\psi_m|\hat P|\psi_n}
= \sum_{\GG, s} c^{s\,*}_m(\GG)\; c^{s}_n(-\GG-\GG_0),
\qquad m,n = 1,\dots,2N,
\end{equation}
and diagonalize it. Because $\hat P$ commutes with $\hat H$ and the set of
$2N$ bands is not cut through the middle of a degenerate level, the
eigenvalues come out as exactly $\pm1$, each in Kramers-degenerate
pairs. If $N_-$ of the $2N$ eigenvalues are $-1$, then
\begin{equation}
\delta_i = (-1)^{N_-/2}.
\label{eq:deltafromcount}
\end{equation}

\begin{algorithm}[H]
\caption{$\mathbb{Z}_2$ invariant of gold from parities}
\label{alg:z2}
\begin{algorithmic}[1]
\State Converge the SCF (SOC on); keep $V_{\mathrm{eff,loc}}$ fixed from
here on (Section~\ref{sec:nscf})
\State $N \gets 6$ \Comment{Kramers pairs below the direct gap}
\State \textbf{Direct-gap check:} \Call{SolveBands}{} on a fine mesh and on
the high-symmetry path; verify
$\min_\kk\,(\varepsilon_{2N+1,\kk} - \varepsilon_{2N,\kk}) > 0$
\For{each of the 8 TRIM $\kk_i = \tfrac12(n_1,n_2,n_3)$}
\State $(\varepsilon, C) \gets$ \Call{SolveBands}{$\kk_i$;\
$N_{\mathrm{want}} = 2N{+}2$;\ $\epsilon = 10^{-9}$}
\State \textbf{assert} $\varepsilon_{2N+1} - \varepsilon_{2N} > 10^{-3}$\,Ha
and bands pair up ($|\varepsilon_{2j} - \varepsilon_{2j-1}| < 10^{-8}$\,Ha)
\State $\mathrm{perm}[g] \gets$ index of plane wave
$(-m_1{-}n_1, -m_2{-}n_2, -m_3{-}n_3)$ for each plane wave $g = (m_1,m_2,m_3)$
\State $\Pi_{mn} \gets \sum_{s,g}\, \overline{C[m,s,g]}\; C[n,s,\mathrm{perm}[g]]$
for $m,n \le 2N$
\Comment{$2N\times2N$, Hermitian}
\State $\lambda \gets \mathrm{eigvalsh}(\Pi)$;\quad \textbf{assert} every
$\big||\lambda_j| - 1\big| < 10^{-6}$
\State $N_- \gets \#\{\lambda_j < 0\}$;\quad \textbf{assert} $N_-$ even;\quad
$\delta_i \gets (-1)^{N_-/2}$
\EndFor
\State $\nu_0 \gets 0$ if $\prod_{i=1}^{8}\delta_i = +1$, else $1$
\State $\nu_j \gets 0$ if $\prod_{i:\,n_j = 1}\delta_i = +1$, else $1$,
for $j = 1,2,3$ \Comment{weak indices: four TRIM each}
\State \Return $(\nu_0;\,\nu_1\nu_2\nu_3)$
\end{algorithmic}
\end{algorithm}

The assertions are not decoration. $|\lambda| \ne 1$ means the band set
cuts a degenerate level, or the eigenvectors are not converged, or the
permutation is wrong; an odd $N_-$ means Kramers pairing is broken.

\section{What to expect for gold, and why}

\begin{center}
\begin{tabular}{@{}lllcc@{}}
\toprule
TRIM & Even pairs (eV) & Odd pairs (eV) & $N_-/2$ & $\delta_i$\\
\midrule
$\Gamma$ ($\times 1$) & $-10.06$, $-5.45\,(\times2)$, $-4.24$, $-3.02\,(\times2)$ & none below $+15.97$ & 0 & $+1$\\
$X$ ($\times 3$) & $-7.32$, $-7.03$, $-2.31$, $-2.09$, $-1.05$ & $+0.71$ & 1 & $-1$\\
$L$ ($\times 4$) & $-7.56$, $-5.57$, $-4.38$, $-2.36$, $-1.63$ & $-1.22$ & 1 & $-1$\\
\bottomrule
\end{tabular}
\end{center}
The table lists the six lowest Kramers pairs at each TRIM with their
energies relative to $E_F$ and their computed parities (reference
calculation: LDA, $E_{\mathrm{cut}}=40$\,Ha). The seventh pair is even at
$X$ ($+5.78$\,eV) and $L$ ($+3.13$\,eV) and odd at $\Gamma$.

\begin{equation}
(-1)^{\nu_0} = \underbrace{(+1)}_{\Gamma}\cdot\underbrace{(-1)^3}_{3X}
\cdot\underbrace{(-1)^4}_{4L} = -1
\qquad\Longrightarrow\qquad \nu_0 = 1 .
\end{equation}
Each weak index involves two $X$ and two $L$ points, giving $+1$: the
result is $(\nu_0;\nu_1\nu_2\nu_3) = (1;000)$. Treat the table as a
prediction to be confirmed by your own numbers, not as an input. The
energies depend on the functional and the cutoff at the level of
$0.1$\,eV; the parities do not.

The pattern can be understood with nothing more than
Sections~\ref{sec:nfe} and \ref{sec:symmetry}:
\begin{itemize}[leftmargin=2em]
\item \textbf{$d$ states are even everywhere.} A Bloch sum of an
atom-centered orbital, evaluated at a TRIM, inherits the parity
$(-1)^\ell$ of the orbital. Five $d$ orbitals give five even pairs at
every TRIM.
\item \textbf{At $\Gamma$} the lowest free-electron state is the constant
plane wave $\GG = 0$, which is even ($s$-like). Odd states are far above
($16$\,eV above $E_F$ in the table).
\item \textbf{At a zone-boundary point} ($X$ or $L$) the two lowest free-
electron states are degenerate, and the crystal potential splits them
into the even, $s$-like $\cos$ state and the odd, $p$-like $\sin$ state
of Eq.~\eqref{eq:nfeparity}. Which one is lower? In the two-plane-wave
model that was decided by the sign of one Fourier component of the
potential. In gold something else dominates: \emph{states of equal parity
repel each other}. The even combination is pushed \emph{up} by the five
even $d$ bands lying just below it; the odd one cannot mix with them and
stays put. In gold (as in copper and silver) this repulsion is strong
enough that the odd state ends up \emph{below} the even one at both $X$
and $L$ --- the inverted order of Section~\ref{sec:nfe}. That is the band
inversion. At $L$ the odd level sits $1.2$\,eV below $E_F$ and the
even one $3.1$\,eV above; the gap between them is the ``$L$ gap''
through which the necks of gold's Fermi surface pass.
\end{itemize}
Because there are four $L$ points, their signs cancel in $\nu_0$: the
strong index is decided by $\Gamma$ versus $X$. The ordering at $L$
matters for \emph{where} on a given surface the protected state appears.

\paragraph{The measurable consequence.} Cut the crystal along a (111)
plane. Photoemission on Au(111) shows a free-electron-like \emph{surface
state} centered at the surface-zone center, about $0.5$\,eV below $E_F$,
sitting inside the bulk $L$ gap. It is split into two parabolas of
opposite spin, shifted in momentum (the Rashba splitting, first resolved
by LaShell, McDougall and Jensen in 1996), with the spin locked
perpendicular to the momentum. Yan \emph{et al.}\ showed that this
long-known ``Shockley state'' is the topological surface state demanded by
$\nu_0 = 1$: the bulk invariant computed here from eight parities is why
that surface band must exist and must be spin--momentum locked. (To see
the surface state itself you would need a slab calculation --- a natural
sequel once the bulk code works.)

\section{Algorithm 2: the invariant from Wilson loops}
\label{sec:wilson}

The parity shortcut exists only for inversion-symmetric crystals. The
general method of Sections~\ref{sec:berry}--\ref{sec:z2def} is worth
implementing: it is what you will need for any crystal without inversion,
it gives an independent check on gold, and it lets you \emph{see} the
topology.

\begin{algorithm}[H]
\caption{Wilson-loop $\mathbb{Z}_2$ on one plane (fixed $k_3 \in
\{0,\tfrac12\}$)}
\label{alg:wilson}
\begin{algorithmic}[1]
\For{$i = 0,\dots,N_2$} \Comment{$k_2 = i/(2N_2)$: from $0$ to $\tfrac12$
inclusive}
\State $(\varepsilon, C_j) \gets$ \Call{SolveBands}{} for the lowest $2N$
bands at $\kk_j = (j/N_1,\, k_2,\, k_3)$, $j = 0,\dots,N_1{-}1$
\State $C_{N_1} \gets C_0$ with every Miller index $m_1$ lowered by one
\Comment{Eq.~\eqref{eq:closure}; do not re-solve}
\State $\mathcal W \gets \mathbf 1_{2N}$
\For{$j = 0,\dots,N_1{-}1$}
\State $M \gets$ overlap matrix \eqref{eq:overlapM} between $C_j$ and
$C_{j+1}$
\State $(U, \Sigma, V^\dagger) \gets \mathrm{svd}(M)$;\quad
record $\min\Sigma$;\quad
$\mathcal W \gets \mathcal W\, U V^\dagger$
\EndFor
\State $\theta[i,\cdot] \gets$ phases of the eigenvalues of $\mathcal W$,
sorted
\EndFor
\State $z \gets$ \Call{CountCrossings}{$\theta$}
\Comment{Algorithm~\ref{alg:count}}
\end{algorithmic}
\end{algorithm}

Line 7 replaces each overlap matrix by its nearest unitary matrix
$UV^\dagger$. On a finite mesh $M$ is not quite unitary (its singular
values are slightly below one, by an amount of order $1/N_1^2$); keeping
only the unitary part makes $\mathcal W$ exactly unitary, so its
eigenvalues lie exactly on the unit circle. A singular value near zero
signals that the group of $2N$ bands is not the same group at the two
neighbouring points --- the direct gap has closed, or the mesh is far too
coarse. To compute the overlaps, scatter both coefficient arrays into the
FFT box (or any common integer index) so that equal $\GG$ line up.

Counting crossings ``by eye'' on a plot is how one checks the result;
a program needs a rule. The standard one (Soluyanov and Vanderbilt)
uses, as the reference line, the middle of the \emph{largest gap} between
adjacent centers, and lets it move with $k_2$: whenever the gap center
jumps from one step to the next, count the centers it jumped over.

\begin{algorithm}[H]
\caption{Counting partner switches by following the largest gap}
\label{alg:count}
\begin{algorithmic}[1]
\Function{CountCrossings}{$\theta[i,j]$,\ $i = 0..N_2$,\ $j = 1..2N$}
\Function{GapCenter}{$\vartheta$} \Comment{$\vartheta$: the $2N$ sorted angles of one row}
\State find adjacent angles (cyclically, the last and the first are
adjacent across $\pm\pi$) with the largest separation;\ \Return their
midpoint on the circle
\EndFunction
\State $z \gets 0$;\quad $g \gets$ \Call{GapCenter}{$\theta[0,\cdot]$}
\For{$i = 1,\dots,N_2$}
\State $g' \gets$ \Call{GapCenter}{$\theta[i,\cdot]$}
\State $z \gets z + \#\{\,j : \theta[i,j] \text{ lies on the shorter arc
between } g \text{ and } g'\,\}$
\State $g \gets g'$
\EndFor
\State \Return $z \bmod 2$
\EndFunction
\end{algorithmic}
\end{algorithm}

Use $N_1 \approx 20$--$40$ points along the loop and enough $k_2$ values
(start with 80) that the curves are smooth; if the count changes when
$N_2$ is doubled, the mesh was too coarse. Run it on the planes $k_3 = 0$
and $k_3 = \tfrac12$ and combine with Eq.~\eqref{eq:nu0fromz}.

\paragraph{What the reference calculation finds.} Figure~\ref{fig:wcc}
shows the twelve charge centers of gold on both planes ($N_1 = 24$,
$N_2 = 80$). Because gold has inversion symmetry the picture is symmetric
under $\theta\to-\theta$ at every $k_2$; most centers stay near
$\theta = 0$ (the $d$ electrons, sitting on the atom) and the action is in
the outermost curves.
\begin{itemize}[leftmargin=2em]
\item \textbf{Plane $k_3 = 0$.} At $k_2 = 0$ the loop passes through
$\Gamma$ and an $L$ point ($\delta_\Gamma\delta_L = -1$): five Kramers
pairs of centers sit at $\theta = 0$ and one at $\theta = \pi$, as
Eq.~\eqref{eq:zakparity} demands. At $k_2 = \tfrac12$ it passes through an
$L$ and an $X$ point ($\delta_L\delta_X = +1$): four pairs at $\theta = 0$
and a mirror-image set at $\theta = \pm0.51\pi$ --- no pair at $\pi$. The
pair that started at $\pi$ has come apart and its members have joined
other partners: $z = 1$.
\item \textbf{Plane $k_3 = \tfrac12$.} Both ends pass through one $L$ and
one $X$ point and look identical (four pairs at $0$, a set at
$\pm0.51\pi$). The curves make large excursions in between but return to
the partners they started with: $z = 0$.
\end{itemize}
Hence $\nu_0 = (1+0) \bmod 2 = 1$ and $\nu_3 = 0$, in agreement with the
parities.

Now look at Fig.~\ref{fig:wcc} critically, because gold is a \emph{hard}
case for this method and the figure shows it. The outer curves do not
move smoothly: near $k_2 \approx 0.26$ and $0.43$ (left) and
$k_2\approx 0.14$ and $0.37$ (right) they appear to jump. True charge
centers are continuous functions of $k_2$ as long as the direct gap is
open; the apparent jumps occur where the loop passes close to the points
at which the gap nearly closes (Section~\ref{sec:directgap}: $0.02$\,eV).
There the wavefunctions of bands 12 and 13 exchange character within a
tiny region of $\kk$-space, the centers sweep across the circle within a
correspondingly tiny interval of $k_2$, and the smallest singular value
of the overlap matrices drops to $0.03$ --- far below the healthy range.
Three consequences:
\begin{itemize}[leftmargin=2em]
\item The \emph{end points} $k_2 = 0, \tfrac12$ are unaffected (the gap
there is large) and, through Eq.~\eqref{eq:zakparity}, already determine
$z$ for an inversion-symmetric crystal. This is the solid part of the
result.
\item A crossing count that runs across an unresolved jump is not
trustworthy. With $N_2 = 80$ Algorithm~\ref{alg:count} returned the
values above; with $N_2 = 20$ it returned $z = 1$ on \emph{both} planes,
i.e.\ the wrong $\nu_0$. An invariant that changes when the mesh is
refined was not converged.
\item The cure is local refinement: wherever the gap center of
Algorithm~\ref{alg:count} moves by more than a fraction of the gap width
between two $k_2$ values, or a singular value falls below $\sim0.3$,
insert intermediate $k_2$ points (and increase $N_1$) until the curves
are continuous. For a material with a comfortable direct gap none of this
is needed; for gold it is the difference between a picture and a proof,
and it is the reason the parity method is the primary tool here.
\end{itemize}
(In an inversion-symmetric crystal the largest gap between centers is
degenerate with its mirror image, so the reference line of
Algorithm~\ref{alg:count} may hop between the two; each hop passes an even
number of centers and does not affect $z$.)

\begin{figure}[H]
\centering
\pgfplotsset{wccplot/.style={width=0.5\textwidth, height=0.42\textwidth,
  xmin=0, xmax=0.5, ymin=-1, ymax=1, xlabel={$k_2$},
  ytick={-1,-0.5,0,0.5,1},
  tick label style={font=\small}, label style={font=\small},
  title style={font=\small}}}
\begin{tikzpicture}
\begin{axis}[wccplot, ylabel={$\theta/\pi$}, title={$k_3 = 0$:\ \ $z = 1$}]
\addplot[only marks, mark=*, mark size=0.6pt, red!70!black] table[x index=0,y index=1] {figdata/wcc_k3_0.dat};
\end{axis}
\end{tikzpicture}\hfill
\begin{tikzpicture}
\begin{axis}[wccplot, yticklabels={}, title={$k_3 = \tfrac12$:\ \ $z = 0$}]
\addplot[only marks, mark=*, mark size=0.6pt, red!70!black] table[x index=0,y index=1] {figdata/wcc_k3_half.dat};
\end{axis}
\end{tikzpicture}
\caption{Wannier charge centers (Wilson-loop phases) of the twelve lowest
spinor bands of gold across half the zone, from the reference calculation
($E_{\mathrm{cut}} = 20$\,Ha, fixed self-consistent potential). $\theta$
is periodic: $+\pi$ and $-\pi$ are the same point. Left: the Kramers pair
at $\theta = \pm\pi$ for $k_2 = 0$ has no counterpart at $k_2 = \tfrac12$
--- partners are switched. Right: the two ends match. The apparent jumps
are unresolved rapid motion near the points of smallest direct gap (see
text).}
\label{fig:wcc}
\end{figure}

\section{Sanity checks}

\begin{enumerate}[leftmargin=2em]
\item \textbf{Empty lattice.} With $V = 0$ at $\Gamma$ the lowest state is
even. At $X$ the two lowest spatial states are degenerate ($\times2$ for
spin): $\Pi$ on those four spinor states must have eigenvalues
$\{+1,+1,-1,-1\}$. This tests the permutation \eqref{eq:parityperm}
without any physics.
\item \textbf{The two-plane-wave model, live.} With the single-cosine
potential of Stage~1 (Chapter~\ref{ch:verify}),
$V(\rr) = 2V_0\cos(\GG_1\cdot\rr)$ and $\kk = \tfrac12\GG_1$, the two
lowest levels must have parities $(+,-)$ in ascending order for $V_0<0$
and $(-,+)$ for $V_0>0$, Eq.~\eqref{eq:nfeparity}. A band inversion made
by hand.
\item \textbf{$\Pi^2 = \mathbf 1$ and $\Pi = \Pi^\dagger$} to rounding,
when the block of bands is not cut through a degenerate level.
\item \textbf{$[\hat P, \hat H] = 0$ and $[\Theta,\hat H]=0$} for random
spinors at each TRIM (Section~\ref{sec:symmetry}).
\item \textbf{Count of even states.} At $\Gamma$ the lowest twelve spinor
bands of gold must all be even.
\item \textbf{SOC off.} The parity counts at the TRIM cannot depend on
SOC: there the groups are separated by $4$--$19$\,eV, far more than any
spin--orbit shift, so setting $k^\ell = 0$ must leave every $\delta_i$
unchanged. What changes is the direct gap away from the TRIM
(Section~\ref{sec:directgap}).
\item \textbf{Parities versus Wilson loops.} On each of the lines
$k_2 \in\{0,\tfrac12\}$ of each plane, the Wilson-loop phases must be
pairwise degenerate, symmetric under $\theta\to-\theta$, and satisfy
Eq.~\eqref{eq:zakparity} with the $\delta_i$ from Algorithm~\ref{alg:z2}.
\item \textbf{Robustness.} Change $E_{\mathrm{cut}}$, the functional, or
the lattice constant by a few percent: every $\delta_i$ and every $z$ must
stay the same. A topological invariant that changes with the cutoff is a
bug.
\end{enumerate}

\begin{gpubox}{GPU/JAX notes}
The parity method is eight small eigenproblems plus one gather
(\texttt{C[..., perm]}) and one $12\times12$ \texttt{eigvalsh} each ---
negligible cost, and it batches over the eight TRIM like any other
$\kk$ axis. The Wilson loop is a fixed-potential band calculation on
$N_1 \times (N_2{+}1)$ points per plane, followed by a \texttt{scan} of
small matrix products; the overlaps are one \texttt{einsum} over the
coefficient arrays once both spheres are scattered onto a common index
set (the FFT box works). The $12\times12$ SVDs and the final
non-Hermitian \texttt{eig} are cheapest on the host. Both methods must run
in double precision with tightly converged eigenvectors.
\end{gpubox}

\begin{keybox}{Summary}
The centers of charge of a group of bands are computable from overlaps of
neighbouring wavefunctions (Wilson loops). With time-reversal symmetry
those centers come in Kramers pairs at the edges of the half zone, and the
$\mathbb{Z}_2$ invariant records whether they switch partners in between.
With inversion symmetry the same bit is the product of the parities of the
Kramers pairs at the eight TRIM, Eq.~\eqref{eq:fukane}, obtained by
applying the permutation $c(\GG) \to c(-\GG-2\kk)$ to your coefficient
arrays. For gold both routes give $\nu_0 = 1$: the odd, $p$-like band lies
below the even, $s$-like band at the zone boundary, and the spin-split
surface state of Au(111) is the visible consequence.
\end{keybox}
